\section{The Yoneda lemma}
We conclude our analysis of the model in $\Cat$ by
providing some further insight into identity elimination.
We saw in \Cref{prop:right-twisted-factorisation-equiv}
that the right twisted arrow factorisation 
\[\Core{A} \to \rTw{A} \to \Core{A} \times A\]
is the comprehensive factorisation of the restricted diagonal
$\Core{A} \to \Core{A} \times A$
into an initial functor followed by a discrete opfibration.
We also saw that it was the algebraic factorisation
of the restricted diagonal into a lari functor
followed by a split opfibration.
It turns out to also be determined by the $\Grpd$-enriched
comprehensive factorisation system on $\Cat$,
which lives ``in between'' the two.
\[ \begin{tikzcd}
	{\text{lari}} & \pitchfork & {\text{split opfibration}} \\
	{1\text{-initial}} & {\bot_{g}} & {\text{groupoidal opfibration}} \\
	{\text{initial}} & \bot & {\text{discrete opfibration}}
	\arrow[from=1-1, to=2-1]
	\arrow[from=2-1, to=3-1]
	\arrow[from=2-3, to=1-3]
	\arrow[from=3-3, to=2-3]
\end{tikzcd} \]
Here we use $\pitchfork$ to indicate that lifts are not unique,
$\bot_g$ when lifts are unique up to unique isomorphism
(unique up to a contractible groupoid),
and $\bot$ when lifts are unique.
The best reference we can find for this result is a comment by Kelly in
\cite[4.7]{kelly1982}.
However, we will only need the following weaker result,
which avoids an explicit description of the 1-initial functors.

\begin{proposition}\label{prop:lari-orthogonal-groupoidal-opfibration}
  Lari functors left-lift against groupoidal opfibrations,
  and those lifts are unique up to unique isomorphism.
\end{proposition}
\begin{proof}
  Consider a diagram of the form
  \[ \begin{tikzcd}
	\AA & \CC \\
	\BB & \DD
	\arrow["G", from=1-1, to=1-2]
	\arrow["L"', from=1-1, to=2-1]
	\arrow["F", from=1-2, to=2-2]
	\arrow["H"', from=2-1, to=2-2]
    \end{tikzcd}\]
  where $(F,l)$ is a groupoidal opfibration and
  $L \vdash R$ is lari.
  We construct a diagonal filler $\phi : \BB \to \CC$.
  Let $b$ be an object in $\BB$.
  Taking the counit of the adjunction $L \vdash R$ at $b$,
  we obtain a lifting problem in $\CC$ over $\DD$.
  \[\begin{tikzcd}
	GRb & {\phi(b)} \\
	HLRb & Hb
	\arrow["l", dashed, from=1-1, to=1-2]
	\arrow["{H \epsilon_b}"', from=2-1, to=2-2]
    \end{tikzcd}\]
  Lifting this using opfibrancy provides an object $\phi(b)$ over $H b$.
  The functorial action of $\phi$ is then determined by
  opcartesianness of the lift $l$.
  The triangle below commutes by definition of $\phi(b)$.
  \[ \begin{tikzcd}
	& \CC \\
	\BB & \DD
	\arrow["F", from=1-2, to=2-2]
	\arrow["\phi", from=2-1, to=1-2]
	\arrow["H"', from=2-1, to=2-2]
    \end{tikzcd}\]  
  To show that the upper triangle commutes,
  \[\begin{tikzcd}
	\AA & \CC \\
	\BB
	\arrow["G", from=1-1, to=1-2]
	\arrow["L"', from=1-1, to=2-1]
	\arrow["\phi"', from=2-1, to=1-2]
    \end{tikzcd}\]
  consider an object $a$ in $\AA$.
  The counit $\epsilon_{La} : L R L a \to L a$
  is the inverse of the unit,
  which is the identity.
  \[\begin{tikzcd}
	{L a} & {L R L a}
	\arrow["{=}", from=1-1, to=1-2]
	\end{tikzcd}\]
  Hence $\epsilon_{La}$ is also the identity.
  Since $(F,l)$ is normal, the lift is also the identity
  \[ \phi L a = GRLa = G a\]
  Hence $\phi$ is a lift for the diagram.
  Any other lift $\phi : \BB \to \CC$ applied to the counit of
  an object $b$ provides a map $\psi(\epsilon_b) : \psi L R b = G R b \to \psi b$
  which also lives over $H \epsilon_b$.
  By opcartesianness of the lift $l$,
  there is a unique morphism $u : \phi b \to \psi b$
  in the fibre over $H b$, which is a groupoid.
  Hence $\phi \iso \psi$ and that isomorphism is unique.
\end{proof}

Now consider a category $A$ with an object $x$.
Taking the pullback of the right twisted arrow category
along a point $x$ in the core,
we obtain the coslice under $x$.
\[
\begin{tikzcd}
	1 & {\Core{A}} \\
	{x /A} & {\rTw{A}} \\
	A & {\Core{A} \times A} \\
	1 & {\Core{A}}
	\arrow["x", from=1-1, to=1-2]
	\arrow["{\rfl x}"', from=1-1, to=2-1]
	\arrow["\rfl", from=1-2, to=2-2]
	\arrow[from=2-1, to=2-2]
	\arrow["\trg"', from=2-1, to=3-1]
	\arrow["\lrcorner"{anchor=center, pos=0.125}, draw=none, from=2-1, to=3-2]
	\arrow["{(\src,\trg)}", from=2-2, to=3-2]
	\arrow[from=3-1, to=3-2]
	\arrow[from=3-1, to=4-1]
	\arrow["\lrcorner"{anchor=center, pos=0.125}, draw=none, from=3-1, to=4-2]
	\arrow["{\pi_1}", from=3-2, to=4-2]
	\arrow["x"', from=4-1, to=4-2]
\end{tikzcd}
\]
Note that since $(\src, \trg) : \rTw A \to (\src, \trg)$
is a discrete opfibration,
so is $\trg : x / A \to A$.
Now suppose we have a groupoidal opfibration $d : C \to A$
with a point $r : 1 \to C$ making the following square commute.
\[ \begin{tikzcd}
	1 & C \\
	{x / A} & A
	\arrow["r", from=1-1, to=1-2]
	\arrow["x"{description}, from=1-1, to=2-2]
	\arrow["{\rfl x}"', from=1-1, to=2-1]
	\arrow["d", from=1-2, to=2-2]
	\arrow["\trg"', from=2-1, to=2-2]
\end{tikzcd} \]
By \Cref{prop:lari-orthogonal-groupoidal-opfibration},
since $\rfl x$ is lari and $d : C \to A$
is a groupoidal opfibration,
there is a diagonal filler $\phi : x / A \to C$.
In fact, since $x / A$ and $C$
are groupoidal fibrations over $A$
(and all functors between groupoidal fibrations
in the slice are morphisms of fibrations)
there is an equivalence of groupoids
\[ \opfibcart(A)(x / A, C) \simeq \{r : 1 \to C \st d r = x\} \]
where both sides have the obvious natural transformations as morphisms.
After straightening (\Cref{thm:opfibration-main}),
the more familiar Yoneda lemma emerges:

\begin{proposition}
  For a category $A$, an object $x \in A$,
  and a functor $C : A \to \Grpd$,
  there is an equivalence of categories
  \[ [A,\Grpd](y (x), C^*) \simeq C^*(x)\]
  where $y(x) : A \to \Grpd$ is given pointwise by the discrete groupoid
  $y(x)(a) := \Hom (x,a)$.
\end{proposition}
\begin{proof}
  Applying \Cref{thm:opfibration-main},
  and noting that $y(x) : A \to \Grpd$ corresponds to the coslice $x / A \in \opfibcart(A)$,
  we have 
  \[ [A,\Grpd](y (x), C^*) \simeq
  \opfibcart(A)(x / A, C) \simeq \{r : 1 \to C \st d r = x\}
  = C^*(x) \]
\end{proof}
